\section{Trace-free Ricci separation from a single monopole}
\label{sec:tracefree-gap}

The parameter estimate of Theorem~\ref{thm:single-parameter} has a
topological consequence at $\beta=8$.  At this value, the scalar and Weyl
terms have exactly the proportions appearing in the Chern--Weil formula
for $2\chi+3\tau$.  We prove the resulting separation estimate for the
trace-free Ricci energy.  Its input is one ordinary monopole on the
negative Yamabe representative of the conformal class.  Two versions
retain the cost of a class whose square does not satisfy the usual
topological bound.

\subsection{Energies and the conformal-class hypothesis}

Let $M$ be a closed, connected, oriented smooth four-manifold.  For a
smooth metric $g$, write
\begin{align}
 V(g)&=\Vol_g(M),&
 S(g)&=\int_M s_g^2\,d\mu_g,\notag\\
 P(g)&=\int_M|\Ric_g^\circ|^2\,d\mu_g,&
 R(g)&=\int_M|\Ric_g|^2\,d\mu_g=P(g)+\frac14S(g).
 \label{eq:trf-energies}
\end{align}
Here $\Ric_g^\circ=\Ric_g-(s_g/4)g$.  Weyl curvature is measured in the
Hilbert--Schmidt norm as an endomorphism of the bundle of two-forms.  Set
\begin{equation}
 T=2\chi(M)+3\tau(M),\qquad
 \E(g)=\inf_{\lambda\in\mathbb R\setminus\{0\}}
 \frac{\int_M|\Ric_g-\lambda g|^2\,d\mu_g}{\lambda^2V(g)},
 \qquad \J(g)=\frac{4P(g)}{R(g)}
 \label{eq:trf-defects}
\end{equation}
whenever $R(g)>0$.

These two scale-invariant defects differ only by scalar variance.  Indeed,
putting $x=1/\lambda$ gives
\[
 \frac{\int_M|\Ric_g-\lambda g|^2\,d\mu_g}{\lambda^2V(g)}
 =\frac{R(g)}{V(g)}x^2
   -\frac{2x}{V(g)}\int_Ms_g\,d\mu_g+4.
\]
Minimizing this quadratic, with the value at $x=0$ understood as an
infimum when necessary, yields
\begin{equation}
 \E(g)=4-\frac{\left(\int_Ms_g\,d\mu_g\right)^2}{V(g)R(g)}
 =\J(g)+\frac{\int_M(s_g-\bar s_g)^2\,d\mu_g}{R(g)},
 \qquad
 \bar s_g=V(g)^{-1}\int_Ms_g\,d\mu_g.
 \label{eq:trf-variance}
\end{equation}
In particular, $0\le\J\le\E\le4$, and $\J=\E$ for a metric of
constant scalar curvature.  If $R(g)=0$, then $g$ is Ricci-flat and the
definition gives $\E(g)=4$; $\J(g)$ is not defined in that case.

For a $\operatorname{Spin}^c$ structure with determinant line bundle $L$,
let $a=c_1(L)\in H^2(M;\mathbb R)$ and
\[
 p=(a_g^+)^2.
\]
The Hodge star on two-forms, and hence the positive harmonic subspace,
is unchanged under a conformal change in dimension four.  Thus $p$ is
constant on $[g]$.

\begin{definition}[Single-monopole hypothesis]
\label{def:trf-input}
For a number $A>0$, a conformal class $[g]$ satisfies the
single-monopole hypothesis if its Yamabe constant is negative and its
representative $\widehat g$ of constant scalar curvature $-4$ admits an
ordinary unperturbed Seiberg--Witten solution in a structure with
$p\ge A$.  We write the solution after a constant spinor rescaling as
\begin{equation}
 D_B\Phi=0,\qquad F_B^+=2i\sigma(\Phi),\qquad
 |\sigma(\Phi)|^2=\frac18|\Phi|^4.
 \label{eq:trf-monopole}
\end{equation}
Thus, if the ordinary normalization is $F_B^+=i\sigma(\Psi)$, then
$\Phi=\Psi/\sqrt2$.  The class $a$, its square, its projection $p$ and
the solution in \eqref{eq:trf-monopole} are fixed throughout the argument
on $[g]$.
\end{definition}

The representative $\widehat g$ is unique with this scalar normalization.
When this hypothesis is imposed on every conformal class of $M$, the
structure may depend on the class.  The number $A$ must be common to
all classes if a uniform conclusion is sought.

\subsection{Scalar curvature and inheritance from conformal solvability}

We first recall the scalar estimate, including its normalization.  This
is the ordinary monopole scalar bound underlying the curvature
inequalities of \cite{LeBrun2001}.

\begin{lemma}[Scalar budget]
\label{lem:trf-scalar}
Suppose that a metric $h$ admits a solution of
\eqref{eq:trf-monopole} with $p=(a_h^+)^2>0$.  With
$s_{h,-}=\max\{-s_h,0\}$, one has
\begin{equation}
 \int_M s_{h,-}^2\,d\mu_h\ge32\pi^2p.
 \label{eq:trf-signed-scalar}
\end{equation}
Consequently, under Definition~\ref{def:trf-input},
\begin{equation}
 16\widehat V\ge32\pi^2p\ge32\pi^2A,
 \qquad \frac{A}{\widehat V}\le\frac1{2\pi^2},
 \qquad \widehat V=V(\widehat g).
 \label{eq:trf-volume-budget}
\end{equation}
\end{lemma}

\begin{proof}
Put $q=|\Phi|^2$.  The integrated spinor Weitzenb\"ock formula in
the normalization \eqref{eq:trf-monopole} is
\[
 0=4\int_M|\nabla_B\Phi|^2\,d\mu_h
       +\int_Ms_hq\,d\mu_h+2\int_Mq^2\,d\mu_h.
\]
It follows that
\[
 2\int_Mq^2\,d\mu_h
 \le\int_Ms_{h,-}q\,d\mu_h
 \le\left(\int_Ms_{h,-}^2\,d\mu_h\right)^{1/2}
      \left(\int_Mq^2\,d\mu_h\right)^{1/2}.
\]
The harmonic projection of $F_B^+$ represents $2\pi i a_h^+$.
Orthogonality of that projection gives
\[
 4\pi^2p\le\int_M|F_B^+|^2\,d\mu_h
             =\frac12\int_Mq^2\,d\mu_h.
\]
In particular $\int q^2>0$.  Squaring the preceding Cauchy--Schwarz
bound and combining the two inequalities proves
\eqref{eq:trf-signed-scalar}.  Applying it to $s_{\widehat g}=-4$
gives \eqref{eq:trf-volume-budget}.
\end{proof}

\begin{proposition}[Inheritance from conformal solvability]
\label{prop:trf-inheritance}
Fix a conformal class $[g]$ and a $\operatorname{Spin}^c$ structure
with $(a_g^+)^2\ge A>0$.  If this same structure has an ordinary
unperturbed solution for every metric in $[g]$, then
Definition~\ref{def:trf-input} holds for $[g]$ and $A$.
\end{proposition}

\begin{proof}
If the Yamabe constant were nonnegative, the class would contain a
representative $h$ of nonnegative constant scalar curvature.  The assumed
solution on $h$ would contradict \eqref{eq:trf-signed-scalar}, since its
left side would vanish and $p>0$.  Equivalently, all terms in the
integrated Weitzenb\"ock formula would be nonnegative, forcing the spinor
and then $F_B^+$ to vanish.  The harmonic projection would then have
$p=0$.

The class therefore has negative Yamabe constant.  Normalize its Yamabe
metric to scalar curvature $-4$.  This is still a metric in $[g]$, so the
same structure supplies an ordinary solution there.  Rescaling the
spinor gives \eqref{eq:trf-monopole}.  The class square and $p$ are
unchanged.
\end{proof}

This proposition allows all-conformal-metric monopole hypotheses to be
used below.  In the proofs that follow, only the resulting solution on
$\widehat g$ is required; no solution is requested at another parameter
value or at another metric in the class.

\subsection{Exact conformal comparison}

The transfer from $\widehat g$ to $g$ retains the trace-free Ricci energy
without introducing a scalar-variance error into the conclusion.

\begin{lemma}[Conformal energy identities]
\label{lem:trf-conformal}
Let $[g]$ have negative Yamabe constant, and write
$g=u^2\widehat g$, where $u>0$ and $s_{\widehat g}=-4$.  Then
\begin{align}
 \delta:=S(g)-S(\widehat g)
 &=48\int_Mu^{-2}|du|_{\widehat g}^2\,d\mu_{\widehat g}
   +36\int_M(u^{-1}\Delta_{\widehat g}u)^2\,d\mu_{\widehat g}
   \ge0,\label{eq:trf-scalar-difference}\\
 P(g)-P(\widehat g)&=\frac{\delta}{12},&
 R(g)-R(\widehat g)&=\frac{\delta}{3},
 \label{eq:trf-conformal-energies}
\end{align}
where $\Delta_{\widehat g}=d^*d$.  In particular, $R(g)>0$, and
\begin{equation}
 \J(g)=
 \frac{R(\widehat g)}{R(\widehat g)+\delta/3}\J(\widehat g)
 +\frac{\delta/3}{R(\widehat g)+\delta/3}.
 \label{eq:trf-convex-comparison}
\end{equation}
If $\J(g)<1$ and $r=P(\widehat g)/\widehat V$, then
\begin{equation}
 \J(\widehat g)\le\J(g)<1,
 \qquad
 r=\frac{4\J(\widehat g)}{4-\J(\widehat g)}
 \le\frac43\J(g).
 \label{eq:trf-small-J}
\end{equation}
\end{lemma}

\begin{proof}
The scalar transformation formula is
\[
 s_g u^3=6\Delta_{\widehat g}u-4u,
 \qquad d\mu_g=u^4d\mu_{\widehat g}.
\]
Squaring and integrating gives
\[
 S(g)=36\int_Mu^{-2}(\Delta_{\widehat g}u)^2\,d\mu_{\widehat g}
       -48\int_Mu^{-1}\Delta_{\widehat g}u\,d\mu_{\widehat g}
       +16\widehat V.
\]
Since $M$ is closed,
\[
 \int_Mu^{-1}\Delta_{\widehat g}u\,d\mu_{\widehat g}
 =\int_M\langle d(u^{-1}),du\rangle_{\widehat g}\,d\mu_{\widehat g}
 =-\int_Mu^{-2}|du|_{\widehat g}^2\,d\mu_{\widehat g}.
\]
Together with $S(\widehat g)=16\widehat V$, this proves
\eqref{eq:trf-scalar-difference}.

In the stated curvature norm, the Chern--Weil identity is
\begin{equation}
 4\pi^2T=\frac1{24}S(h)
       +2\int_M|W_h^+|^2\,d\mu_h-\frac12P(h)
 \label{eq:trf-chern-weil}
\end{equation}
for every metric $h$.  The Weyl integral is conformally invariant.
Subtracting \eqref{eq:trf-chern-weil} for $g$ and $\widehat g$ gives
$P(g)-P(\widehat g)=\delta/12$.  The identity
$R=P+S/4$ then gives the second formula in
\eqref{eq:trf-conformal-energies}.

Now $R(\widehat g)=P(\widehat g)+4\widehat V>0$, and substitution
in $\J=4P/R$ proves \eqref{eq:trf-convex-comparison}.  This expresses
$\J(g)$ as a convex combination of $\J(\widehat g)$ and $1$, with
a strictly positive coefficient on $\J(\widehat g)$.  If $\J(g)<1$,
it follows that $\J(\widehat g)\le\J(g)<1$.  Finally
$\J(\widehat g)=4r/(r+4)$, so
\[
 r=\frac{4\J(\widehat g)}{4-\J(\widehat g)}
 \le\frac43\J(\widehat g)\le\frac43\J(g).
\]
\end{proof}

\subsection{The topological calibration and the separation theorem}

For the fixed class $a$, define
\begin{equation}
 D_\beta(h,a)=
 (1+\beta^{-1})S(h)+6(\beta+1)\int_M|W_h^+|^2\,d\mu_h
 -72\pi^2p,
 \qquad \beta\ge2.
 \label{eq:trf-pencil}
\end{equation}
For the rest of this subsection, all terms in $D_\beta$, $r$ and $b$
are evaluated at $\widehat g$, and
\[
 r=\frac{P(\widehat g)}{\widehat V},\qquad
 b=\frac{(-a^2)_+}{\widehat V}.
\]
Write $c_{\mathrm{par}}>0$ for a universal constant furnished by
Theorem~\ref{thm:single-parameter}.  Its single-solution estimate is
\begin{equation}
 (1+b)\left(r+\frac{D_\beta(\widehat g,a)}{\beta\widehat V}\right)
 \ge c_{\mathrm{par}}\frac{(\beta-2)^2}{\beta^4}.
 \label{eq:trf-parameter}
\end{equation}
In particular, with
\begin{equation}
 H=r+\frac{D_8(\widehat g,a)}{8\widehat V},\qquad
 c_8=\frac{9c_{\mathrm{par}}}{1024},\qquad
 \gamma=\frac{4c_8}{27}=\frac{c_{\mathrm{par}}}{768},
 \label{eq:trf-constants}
\end{equation}
we have $(1+b)H\ge c_8>0$.

The ratio of the Weyl and scalar coefficients in
\eqref{eq:trf-pencil} is $6\beta$.  It equals the corresponding ratio
$48$ in \eqref{eq:trf-chern-weil} precisely when $\beta=8$.
Multiplication of \eqref{eq:trf-chern-weil} by $27$ therefore gives
the exact identity
\begin{equation}
 D_8(\widehat g,a)
 =108\pi^2\left(T-\frac23p\right)+\frac{27}{2}P(\widehat g).
 \label{eq:trf-beta-eight}
\end{equation}
Set
\begin{equation}
 d=\left(\frac{T}{A}-\frac23\right)_+.
 \label{eq:trf-topological-excess}
\end{equation}
Since $p\ge A$, one has $T-2p/3\le Ad$.  Lemma~\ref{lem:trf-scalar}
and \eqref{eq:trf-beta-eight} imply
\begin{equation}
 H\le\frac{43}{16}r+\frac{27}{4}d.
 \label{eq:trf-H-upper}
\end{equation}
If $\J(g)<1$, Lemma~\ref{lem:trf-conformal} consequently gives
\begin{equation}
 H\le\frac{43}{12}\J(g)+\frac{27}{4}d
       \le\frac{27}{4}\bigl(\J(g)+d\bigr).
 \label{eq:trf-H-J}
\end{equation}
The numerical inequality in the last step is $43/12\le27/4$.

\begin{theorem}[Trace-free Ricci separation]
\label{thm:trf-gap}
Suppose that $[g]$ satisfies Definition~\ref{def:trf-input} and that its
selected class satisfies $a^2\ge T$.  There is a universal
$\varepsilon_0>0$ such that
\begin{equation}
 \boxed{\J(g)+\left(\frac{T}{A}-\frac23\right)_+
             \ge\varepsilon_0.}
 \label{eq:trf-gap}
\end{equation}
For the constant in \eqref{eq:trf-parameter}, one may take
\begin{equation}
 K_\pi=1+\frac1{6\pi^2},\qquad
 \varepsilon_0=\min\left\{1,\frac{\gamma}{K_\pi}\right\}.
 \label{eq:trf-epsilon}
\end{equation}
The conclusion also holds with $\E(g)$ in place of $\J(g)$.
\end{theorem}

\begin{proof}
From \eqref{eq:trf-chern-weil},
\begin{equation}
 4\pi^2T\ge-\frac12P(\widehat g),\qquad
 \frac{(-T)_+}{\widehat V}\le\frac{r}{8\pi^2}.
 \label{eq:trf-negative-T}
\end{equation}
The assumption $a^2\ge T$ thus gives $b\le r/(8\pi^2)$.
If $\J(g)\ge1$, \eqref{eq:trf-gap} follows from
$\varepsilon_0\le1$.  Otherwise \eqref{eq:trf-small-J} gives
$r<4/3$ and hence $1+b\le K_\pi$.  As $H>0$, the parameter estimate
and \eqref{eq:trf-H-J} imply
\[
 c_8\le(1+b)H
 \le K_\pi\frac{27}{4}\bigl(\J(g)+d\bigr).
\]
This proves \eqref{eq:trf-gap} with \eqref{eq:trf-epsilon}.
Finally \eqref{eq:trf-variance} gives $\E(g)\ge\J(g)$.
\end{proof}

The constants in this proof are obtained from the parameter estimate
with its factor $1+b$ retained.  In particular, the absorption of that
factor is accounted for by $K_\pi$; no equality between constants in
the estimates before and after absorption is needed.

\subsection{Two class-square costs}

The condition $a^2\ge T$ can be replaced by a quantitative loss.  There
are two useful ways to measure it:
\begin{equation}
 Y=(T-a^2)_+,\qquad y=\frac{Y}{A},\qquad
 \nu=\frac{(-a^2)_+}{A}.
 \label{eq:trf-cost-definitions}
\end{equation}
The first compares the class square with the topological number $T$;
the second records only its negative part.

\begin{corollary}[Negative class square]
\label{cor:trf-nu}
Under Definition~\ref{def:trf-input}, with no restriction on $a^2$,
\begin{equation}
 \boxed{\J(g)+d\ge\frac{c_*}{1+\nu/(2\pi^2)},
 \qquad c_*=\min\{1,\gamma\}.}
 \label{eq:trf-nu-cost}
\end{equation}
In particular, $a^2\ge0$ implies $\J(g)+d\ge c_*$, without a
further comparison between $a^2$ and $T$.
\end{corollary}

\begin{proof}
By \eqref{eq:trf-volume-budget},
\[
 b=\nu\frac{A}{\widehat V}\le\frac{\nu}{2\pi^2}.
\]
If $\J(g)<1$, then $H>0$ and \eqref{eq:trf-H-J} give
\[
 c_8\le(1+b)H
 \le\left(1+\frac{\nu}{2\pi^2}\right)
       \frac{27}{4}\bigl(\J(g)+d\bigr).
\]
This proves the result with numerator $\gamma$ on this branch.  If
$\J(g)\ge1$, the right side of \eqref{eq:trf-nu-cost} is at most
$1$, so the choice $c_*\le1$ proves the other branch.
\end{proof}

\begin{corollary}[Square deficit relative to topology]
\label{cor:trf-y}
Under Definition~\ref{def:trf-input}, with no restriction on $a^2$,
\begin{equation}
 \boxed{\J(g)+d\ge\frac{c_y}{1+y},
 \qquad c_y=\min\left\{1,\frac{\gamma}{K_\pi}\right\}.}
 \label{eq:trf-y-cost}
\end{equation}
Thus a bound $y\le L$ gives $\J(g)+d\ge c_y/(1+L)$.
\end{corollary}

\begin{proof}
The elementary inequality $(-a^2)_+\le(-T)_++Y$, together with
\eqref{eq:trf-negative-T} and \eqref{eq:trf-volume-budget}, gives
\[
 b\le\frac{r}{8\pi^2}+\frac{y}{2\pi^2}.
\]
When $\J(g)<1$, it follows that
\[
 1+b\le K_\pi+\frac{y}{2\pi^2}\le K_\pi(1+y).
\]
Here the last inequality uses $1/(2\pi^2)\le K_\pi$.  As before,
\eqref{eq:trf-H-J} and $(1+b)H\ge c_8$ imply
\begin{equation}
 \left(K_\pi+\frac{y}{2\pi^2}\right)
          \bigl(\J(g)+d\bigr)\ge\gamma
 \qquad\text{if }\J(g)<1.
 \label{eq:trf-y-refined-branch}
\end{equation}
Replacing the first factor by $K_\pi(1+y)$ yields the asserted
inequality on this branch.  The branch $\J(g)\ge1$ is covered by
$c_y\le1$.
\end{proof}

\begin{remark}
The two costs are complementary.  If $0\le a^2<T$, then $y>0$ but
$\nu=0$, and Corollary~\ref{cor:trf-nu} has no class-square loss.
If $T\le a^2<0$, then $y=0$ but $\nu>0$, and
Theorem~\ref{thm:trf-gap} absorbs the negative square using topology.
In particular, $\nu\le y$ is not a valid comparison without a further
sign condition on $T$.  Neither corollary asserts an optimal dependence
on its cost, and neither permits the constant to be kept fixed while
the corresponding cost tends to infinity.
\end{remark}

\subsection{The region with a positive Ricci gap}

The preceding inequalities yield a positive lower bound for $\J$ only
after the topological excess is subtracted.  For the selected class,
they give
\begin{equation}
 \J(g)\ge\frac{c_y}{1+y}-d,
 \qquad
 \J(g)\ge\frac{c_*}{1+\nu/(2\pi^2)}-d.
 \label{eq:trf-positive-region}
\end{equation}
The corresponding conclusion is nontrivial when its right side is
strictly positive.  A bounded square cost alone does not impose that
condition on $T/A$.

\begin{corollary}[Uniform separation and Einstein obstruction]
\label{cor:trf-uniform}
Suppose that every conformal class of $M$ satisfies
Definition~\ref{def:trf-input} with the same $A>0$.  It is enough that
the selected structures satisfy one of the following conditions:
\begin{enumerate}
 \item $a^2\ge T$ in every class, and $d<\varepsilon_0$;
 \item $y\le L$ in every class for a fixed $L\ge0$, and
       $d<c_y/(1+L)$;
 \item $\nu\le N$ in every class for a fixed $N\ge0$, and
       $d<c_*/(1+N/(2\pi^2))$.
\end{enumerate}
In the respective cases, every smooth metric on $M$ satisfies
$\J(g)\ge\eta$ and $\E(g)\ge\eta$, where
\[
 \eta=\varepsilon_0-d,\qquad
 \eta=\frac{c_y}{1+L}-d,\qquad\text{or}\qquad
 \eta=\frac{c_*}{1+N/(2\pi^2)}-d.
\]
In each case $\eta>0$, and $M$ admits no Einstein metric.
\end{corollary}

\begin{proof}
Apply Theorem~\ref{thm:trf-gap}, Corollary~\ref{cor:trf-y}, or
Corollary~\ref{cor:trf-nu} in each conformal class.  All quantities on
the displayed right sides are uniform over metrics on the fixed
manifold.  The inequality for $\E$ follows from
\eqref{eq:trf-variance}.

An Einstein metric with nonzero Einstein constant has $R>0$ and $P=0$,
contradicting $\J\ge\eta>0$.  A Ricci-flat metric requires a separate
argument: its conformal class has Yamabe constant zero, contrary to the
assumption that every conformal class is negative.  Under the stronger
conformal solvability hypothesis of Proposition~\ref{prop:trf-inheritance},
the same exclusion follows directly from
$\int_Ms_-^2\,d\mu\ge32\pi^2A>0$.
\end{proof}

For example, whenever an actual application has
$T=A-k-4\ell$ and $a^2\ge T$, \eqref{eq:trf-gap} becomes
\[
 \J(g)+\left(\frac13-\frac{k+4\ell}{A}\right)_+
 \ge\varepsilon_0.
\]
The positive budget and the actual monopole hypothesis remain necessary
inputs to this substitution.  No conclusion at $A=0$ is obtained by
continuity from the normalized statement.

The same uniform gap obstructs pointwise approximation.  If
\[
 \delta_\sigma(g)=
 \|g^{-1}\Ric_g-\sigma I\|_{L^\infty,\mathrm{op}},
 \qquad \sigma\in\{-1,1\},
\]
then testing the definition of $\E$ with $\lambda=\sigma$ gives
$\E(g)\le4\delta_\sigma(g)^2$.  Thus the conclusion of
Corollary~\ref{cor:trf-uniform} implies
$\delta_\sigma(g)\ge\sqrt\eta/2$ for either sign.

\subsection{Finite covers and twisted monopoles}

The argument also applies when the actual ordinary spinor lives on a
finite cover.  The precise statement is useful because the covering
manifold need not supply monopoles on arbitrary metrics.

\begin{proposition}[Finite-cover compatibility]
\label{prop:trf-cover}
Let $\pi:\widetilde M\to M$ be a finite connected oriented cover of degree $m$,
and let $[g]$ have negative Yamabe constant.  Suppose that
$\pi^*\widehat g$ admits an ordinary solution with determinant class
$\widetilde a$ and
\[
 \widetilde p=(\widetilde a_{\pi^*\widehat g}^+)^2\ge mA,
 \qquad A>0.
\]
Put
\[
 y_\pi=\frac{(mT-\widetilde a^2)_+}{mA},\qquad
 \nu_\pi=\frac{(-\widetilde a^2)_+}{mA}.
\]
Then \eqref{eq:trf-nu-cost} and \eqref{eq:trf-y-cost} hold on $M$
with $\nu_\pi,y_\pi$ in place of $\nu,y$.  If
$\widetilde a^2\ge mT$, then \eqref{eq:trf-gap} holds on $M$.
The constants are independent of $m$.
\end{proposition}

\begin{proof}
The pullback $\pi^*\widehat g$ has scalar curvature $-4$.  It is the
normalized negative Yamabe representative of $[\pi^*g]$.  Pullback
and integration over the covering sheets give
\begin{gather*}
 \widetilde V=mV,\qquad
 \widetilde S=mS,\qquad
 \widetilde P=mP,\qquad
 \widetilde R=mR,\qquad
 \widetilde T=mT,\\
 \J(\pi^*g)=\J(g),\qquad \E(\pi^*g)=\E(g).
\end{gather*}
The last equality also follows directly from the defining quotient for
each $\lambda\ne0$.  Apply the results above on $\widetilde M$ with
budget $\widetilde A=mA$.  Its topological excess is
$(\widetilde T/\widetilde A-2/3)_+=d$, and its two normalized costs
are exactly $y_\pi$ and $\nu_\pi$.  The resulting inequalities are
therefore the asserted inequalities downstairs, with the same constants.
\end{proof}

For a regular cover, the normalized negative Yamabe representative is
deck-invariant by uniqueness.  To obtain the input from conformal
solvability, fix one ordinary structure upstairs with
$\widetilde p\ge mA>0$ on the lifted conformal class.  Suppose that it
has solutions for every lifted metric in that class.  A nonnegative
Yamabe representative downstairs would lift to a nonnegative
scalar-curvature metric and contradict Lemma~\ref{lem:trf-scalar}
upstairs.  Thus the base class is negative, and its normalized
representative can be lifted and used in Proposition~\ref{prop:trf-cover}.
Only these lifted metrics are involved.  The proof makes no solvability
assumption for non-invariant metrics, and it does not require the
ordinary solution itself to be deck-invariant.

For the twisted case, let $\pi:\widetilde M\to M$ be a nontrivial
double cover with associated real sign local system $\ell$.  Suppose
that a genuine $\operatorname{Spin}^{c_-}$ structure has twisted Euler
class $a\in H^2(M;\ell)$ with $p=(a_g^+)^2\ge A>0$, and an actual unperturbed
$\operatorname{Pin}^-(2)$ monopole on the negative Yamabe
representative.  The lifting theorem of \cite{IMN2016} gives an
ordinary $\operatorname{Spin}^c$ solution upstairs, after matching the
constant spinor normalization, with $\widetilde a=\pi^*a$.
Since pullback commutes with the twisted Laplacian and the Hodge star,
\begin{equation}
 \widetilde a^2=2a^2,
 \qquad \widetilde p=2(a_g^+)^2,
 \qquad \widetilde T=2T.
 \label{eq:trf-pin-scaling}
\end{equation}
Products of two sign-local-system classes take values in the trivial
local system, so these squares and the associated integrals are defined
downstairs.  With $\widetilde A=2A$, the quantities $T/A$, $y$ and
$\nu$ are unchanged.  Proposition~\ref{prop:trf-cover} therefore
proves the same separation and cost inequalities for these twisted
budgets, without a covering-degree loss.

If the twisted solution is instead supplied throughout the base
conformal class, the preceding inheritance argument first supplies its
negative Yamabe representative.  Neither formulation asks for twisted
solvability on every conformal class unless an all-metric conclusion on
$M$ is intended.  In that conclusion the positive projection budget and
any bound on $y$ or $\nu$ must be verified for the actual selected
twisted structures in each class.  In particular, the square condition
$a^2\ge T$ is a separate hypothesis, not a consequence of the existence
of an arbitrary twisted class.
